%% machines-effective-procedures-supplement.tex --- separately authored additions
\section*{Editorial supplement to Machines and Effective Procedures}
\addcontentsline{toc}{section}{Editorial supplement to Machines and Effective Procedures}

\noindent\textit{Reader note.} The following separately authored material
retains the nonoverlapping recursive-function and machine-model additions from
the preserved historical draft.  Consumer references now point to the exact
inherited definitions that precede it.  The supplement neither replaces nor
alters the exact \emph{Machines and Effective Procedures} chapter
\cite{Kleene1952,ShepherdsonSturgis1963}.

%% Necessity: the exact inherited chapter does not contain the recursive-
%% function equivalence, URM formalism, or Church--Turing boundary required here.
%% Definition: this supplement is an independently revised presentation of the
%% distinct topics preserved in historical body lines 221--400, checked against
%% the displayed primary and modern sources.
%% Standing: it is editorial handbook content, not inherited restoration or a
%% claim that the historical derivative remains byte exact.
%% Nonredundancy: the hash-pinned historical original remains unreachable, while
%% the exact inherited chapter supplies the machine core without these topics.
%% Use: the Part IV driver places this supplement immediately after that exact
%% chapter and before the universal-computation chapter it references.
%% BEGIN EDITORIAL P4 MACHINES SUPPLEMENT
\section[Minimization and Formalisms]{Unbounded minimization and equivalent formalisms}
\label{sec:p04-partial-recursion}

\begin{hbdefinition}{def:p04-partial-recursive}{Unbounded minimization and partial recursive functions}
For a partial map \(g:\N^{k+1}\rightharpoonup\N\), its \emph{unbounded
minimization} is the partial map
\[
  (\mu g)(\vec x)=\mu y\,[g(\vec x,y)=0].
\]
This value is defined as the least \(y\) for which \(g(\vec x,y)\) is defined
and equals \(0\), provided \(g(\vec x,z)\) is defined and nonzero for every
\(z<y\); it is undefined if an earlier required computation diverges or no
such zero is reached.  For a total predicate \(R(\vec x,y)\), the relational
notation
\[
 \mu y\,R(\vec x,y)
 \defeq \mu y\,[z_R(\vec x,y)=0],
 \qquad
 z_R(\vec x,y)=
 \begin{cases}0,&R(\vec x,y),\\1,&\neg R(\vec x,y),\end{cases}
\]
is an eliminable least-true abbreviation.  It is distinct in type from the
least-zero operator on partial numerical maps and is used in the normal-form
display below.

For each positive arity, the \emph{partial recursive functions}, also called the
\emph{\(\mu\)-recursive functions}, are the least class of
partial maps \(\N^k\rightharpoonup\N\) with \(k\geq1\) that contains the zero,
successor, and projection functions and is closed under composition, primitive
recursion, and unbounded minimization.  Composition is defined only when all
required inner values and then the outer value are defined; a recursively
defined value is defined only when its base value and every preceding value
needed by the step clause are defined.  This is a class of partial numerical
functions, not a class of formulas or machine descriptions.  It strictly adds
the possibility of divergence to the total class in
the named primitive-recursive-functions definition in
Section~\ref{sec:turing-machines-02}; that distinction is
needed by the normal-form and machine-equivalence theorems below.
\index{unbounded minimization}
\index{partial recursive function}
\index{mu-recursive function@\(\mu\)-recursive function}
\end{hbdefinition}

\begin{hbtheorem}{thm:p04-kleene-normal-form}{cited=Kleene1952}{Kleene normal form}
For each \(k\ge1\), there are a primitive recursive relation
\(\mathsf T_k(e,x_1,\ldots,x_k,y)\) and a primitive recursive function
\(\mathsf U(y)\) such that every partial recursive
\(f:\N^k\rightharpoonup\N\) has an index \(e\), in the fixed computation
coding of this normal-form theorem, satisfying
\[
  \operatorname{dom}(f)
  =\{\bar x:\exists y\,\mathsf T_k(e,\bar x,y)\},
  \qquad
  f(\bar x)=\mathsf U\bigl(\mu y\,
    \mathsf T_k(e,\bar x,y)\bigr)
  \quad\text{for }\bar x\in\operatorname{dom}(f),
\]
where the minimization selects the least \(y\) satisfying \(\mathsf T_k\).
Conversely, every partial function displayed in this form is partial
recursive.  Here \(\mathsf T_k\) is the relation saying that \(y\) is a valid
finite accepting-computation code under the computation coding fixed for this
normal form, and \(\mathsf U\) extracts its numerical output; it is not the
bounded object-language computation formula introduced later for
arithmetization.
\end{hbtheorem}

\begin{hbtheorem}{thm:p04-tm-partial-recursive}{cited=Kleene1952}{Turing-computable iff partial recursive}
For every \(k\ge1\) and every partial map
\(f:\N^k\rightharpoonup\N\), the following are equivalent:
\begin{enumerate}
  \item \(f\) is computed by a Turing machine under
        the named Turing-machine definition in
        Section~\ref{sec:turing-machines-01} and
        Definition~\ref{def:partial-computable};
  \item \(f\) is partial recursive in the closure sense of
        Definition~\ref{def:p04-partial-recursive}.
\end{enumerate}
\end{hbtheorem}

The equivalence is at the level of partial numerical functions, after the
input/output encodings are fixed.  Historical equivalences among general recursion,
Turing computation, and lambda definability are in
\cite{Kleene1936,Turing1937Lambda}; \cite{Kleene1952} gives the modern
partial-function development.

The index \(e\) in Theorem~\ref{thm:p04-kleene-normal-form} belongs to that
cited normal-form coding; it is not yet being identified with an index in the
handbook's bespoke program numbering.  The next chapter fixes the single
effective program numbering used by every index-dependent construction in this
part and verifies its enumeration and parameter-specialization properties.  No
second handbook program numbering or unrelated compiler is introduced.

\begin{hbdefinition}{def:p04-shepherdson-sturgis-urm}{Shepherdson--Sturgis Unlimited Register Machine}
The \emph{Shepherdson--Sturgis Unlimited Register Machine} (URM) has registers
\(1,2,\ldots\), each storing a natural number, and a finite program mentioning
only finitely many registers.  A program has lines \(1,\ldots,\ell\).  Each
non-normal exit placeholder \(E_j\) in an instruction is replaced on that
line by a target in \(\{1,\ldots,\ell+1\}\), where line \(\ell+1\) is the
halting exit; the normal exit from line \(i\) is line \(i+1\).

Its six instruction forms are \(P(n)\), increment register \(n\);
\(D(n)\), decrement register \(n\) when its value is positive;
\(O(n)\), set register \(n\) to zero; \(C(m,n)\), copy register \(m\) to
register \(n\), with \(m\ne n\); \(J[E_1]\), take the displayed exit
unconditionally; and \(J(m)[E_1]\), take exit \(E_1\) exactly when register
\(m\) is zero and otherwise take the normal exit.  The four data instructions
\(P,D,O,C\) take the normal exit.  In accordance with the primary URM
presentation, \(D(n)\) has no transition at zero; a run that attempts it is
undefined and supplies no computed value.  All unmentioned registers are
unchanged.  Execution begins at line \(1\) and halts only on reaching line
\(\ell+1\).

To compute \(f:\N^n\rightharpoonup\N\), choose distinct register addresses
\(x_1,\ldots,x_n,y\) and a bound \(N\) containing them.  Initialize the
\(x_i\)-registers with the inputs.  If \(f(\vec x)\) is undefined, the program
has no halting-exit outcome; if it is defined, the program reaches the halting
exit with register \(y\) containing \(f(\vec x)\), while registers
\(1,\ldots,N\) other than \(y\) retain their initial contents.  Registers
above \(N\) may be scratch space.
This name denotes that exact formalism and function convention, not an
unspecified generic register machine; the distinction is required for its row in
Table~\ref{tab:p04-machine-equivalences} \cite{ShepherdsonSturgis1963}.
\index{Unlimited Register Machine!Shepherdson--Sturgis}
\index{register machine!Shepherdson--Sturgis}
\end{hbdefinition}

\begin{landscape}
\begingroup\small\setlength{\tabcolsep}{3pt}
\begin{longtable}{@{}p{1.10in}p{2.25in}p{1.85in}p{1.84in}@{}}
\caption{Machine models, recursive-function formalisms, and their exact
evidence relationships.}
\label{tab:p04-machine-equivalences}\\
\toprule
Model & Object and function convention & Verified relationship & Evidence \\
\midrule
\endfirsthead
\toprule
Model & Object and function convention & Verified relationship & Evidence \\
\midrule
\endhead
Fixed single-tape Turing machine & Partial numerical functions; convergence
means accepting with a valid canonical output, while rejection, divergence,
and invalid acceptance are undefined. & Defines the Turing-computable class
under exactly this convergence convention. &
the named Turing-machine definition in Section~\ref{sec:turing-machines-01}
and Definition~\ref{def:partial-computable}; the next chapter supplies the index
realization. \\
\addlinespace
Partial recursive functions & Least partial-function class containing zero,
successor, and projections and closed under composition, primitive recursion,
and the definedness-sensitive \(\mu\)-operator. & Same partial numerical
functions as the fixed Turing-machine formalism. & Cited
Theorem~\ref{thm:p04-tm-partial-recursive};
\cite{Kleene1936,Kleene1952,Turing1937Lambda}. \\
\addlinespace
Shepherdson--Sturgis URM & Countably many \(\N\)-valued registers and finite
programs using exactly \(P,D,O,C,J,J\)-zero as specified in
Definition~\ref{def:p04-shepherdson-sturgis-urm}. & Computes exactly the
partial recursive numerical functions; equivalence with the preceding Turing
model follows through that shared class. & Primary URM construction and
two-symbol-machine reduction: \cite{ShepherdsonSturgis1963}; the final
TM--URM edge is transitive. \\
\addlinespace
Lambda definability & Fix Church numerals; on represented inputs, a term
represents a numerical function when its application to the input numerals
converts to the numeral for the value.  Comparison is by represented
functions, not raw-term equality. & Turing's historical theorem identifies the
corresponding computable, lambda-definable, and general recursive function
classes on positive integers. & \cite{Turing1937Lambda}; no modern
undefined-input convention is attributed to that paper. \\
\addlinespace
Effective calculability & Informal explicandum, not a fourth formal model or a
class fixed by a formal definition. & Church--Turing thesis boundary; no
equivalence edge is asserted from a theorem. &
\cite{Church1936,Turing1937Lambda,Soare1996ComputabilityRecursion}; see
Boundary~\ref{def:p04-church-turing}. \\
\bottomrule
\end{longtable}
\endgroup
\end{landscape}

\begin{proofboundary}{def:p04-church-turing}{Church--Turing thesis}
The \emph{Church--Turing thesis} identifies the informally understood
effectively calculable numerical functions with the Turing-computable
functions.  Church proposed identifying effective calculability with
lambda-definability---equivalently, by Kleene's theorem, general
recursiveness---while Turing analyzed machine computation.  The equivalence
and robustness of the resulting formal models support the identification
\cite{Church1936,Kleene1936,Turing1937Lambda,
Soare1996ComputabilityRecursion}.  They do
not prove it, because ``effective calculability'' is the informal explicandum,
not an independently formalized class quantified over by a theorem.  The
thesis stated here is not the claim that every physical process is Turing
computable, and formal-model equivalence does not imply equality of complexity
costs (the named formalism-versus-complexity warning in
Section~\ref{sec:turing-machines-01}).
\index{Church--Turing thesis}
\end{proofboundary}
%% END EDITORIAL P4 MACHINES SUPPLEMENT
